%%
%% preamble.tex — Shared LaTeX preamble for all papers in the folio.
%%
%% This file is a TEMPLATE.  Paper-specific content (title, authors,
%% macros, abstract, chapter \input lines) is generated by the content
%% pipeline from each paper's .ts manifest.
%%
%% Do NOT put paper-specific definitions here.
%%

\documentclass[12pt,a4paper]{report}

% ── Packages ──────────────────────────────────────────────────────────────────
\usepackage[utf8]{inputenc}
\usepackage[T1]{fontenc}
\usepackage{amsmath}
\usepackage{amssymb}
\usepackage{amsthm}
\usepackage{mathtools}
\usepackage{hyperref}
\usepackage{geometry}
\usepackage{graphicx}
\usepackage{cite}
\usepackage{ccicons}          % CC icon glyphs
\usepackage{titlesec}         % Compact heading styles
\usepackage{tikz}             % Inline figures
\usepackage{tikz-cd}          % Commutative diagrams (tikzcd environment)
\usepackage[dvipsnames]{xcolor}  % Extended colour names (teal, etc.)
\usepackage{booktabs}            % Professional table rules (\toprule, \midrule, \bottomrule)
\usepackage{enumitem}            % Customised list labels (label=, leftmargin=)
\usepackage{mathrsfs}            % \mathscr script letters
\usepackage{textcomp}            % \textdegree, \textbullet, etc.
\usepackage{newunicodechar}      % Map Unicode glyphs to LaTeX commands

% ── Unicode character mappings ───────────────────────────────────────────────
% Allow Unicode symbols in .md content to pass through pdflatex.
\newcommand{\unimathsafe}[1]{\ifmmode #1\else $#1$\fi}
\newunicodechar{✓}{\unimathsafe{\checkmark}}
\newunicodechar{✗}{\unimathsafe{\times}}
\newunicodechar{✔}{\unimathsafe{\checkmark}}
\newunicodechar{★}{\unimathsafe{\bigstar}}
\newunicodechar{°}{\ifmmode ^{\circ}\else $^{\circ}$\fi}
\newunicodechar{×}{\unimathsafe{\times}}
\newunicodechar{←}{\unimathsafe{\leftarrow}}
\newunicodechar{→}{\unimathsafe{\rightarrow}}
\newunicodechar{≤}{\unimathsafe{\leq}}
\newunicodechar{≥}{\unimathsafe{\geq}}
\newunicodechar{≈}{\unimathsafe{\approx}}
\newunicodechar{∞}{\unimathsafe{\infty}}
\newunicodechar{Δ}{\unimathsafe{\Delta}}
\newunicodechar{β}{\unimathsafe{\beta}}
\newunicodechar{⊙}{\unimathsafe{\odot}}
\newunicodechar{●}{\textbullet}
\newunicodechar{○}{$\circ$}
\newunicodechar{…}{\ldots}
\newunicodechar{⚠}{\textbf{!}}
% Subscript digits (U+2080–U+2089)
\newunicodechar{₀}{\textsubscript{0}}
\newunicodechar{₁}{\textsubscript{1}}
\newunicodechar{₂}{\textsubscript{2}}
\newunicodechar{₃}{\textsubscript{3}}
\newunicodechar{₄}{\textsubscript{4}}
\newunicodechar{₅}{\textsubscript{5}}
\newunicodechar{₆}{\textsubscript{6}}
\newunicodechar{₇}{\textsubscript{7}}
\newunicodechar{₈}{\textsubscript{8}}
\newunicodechar{₉}{\textsubscript{9}}

% ── Custom symbols ───────────────────────────────────────────────────────────
\newcommand{\bigbowtie}{\mathop{\triangleright\!\!\!\triangleleft}}

\geometry{margin=1in}

% ── Suppress hyperref math-in-bookmark warnings ─────────────────────────────
% Math in theorem titles generates "Token not allowed in PDF string" warnings.
% These are harmless (hyperref strips the math for bookmarks) but noisy.
\hypersetup{pdfencoding=auto, psdextra}
\pdfstringdefDisableCommands{%
  \let\mathbf\relax
  \let\mathrm\relax
  \let\mathcal\relax
  \let\mathfrak\relax
  \let\mathbb\relax
  \let\Sigma\relax
  \let\Delta\relax
  \let\not\relax
}

% ── Compact chapter/section formatting (no page breaks, normal-size fonts) ────
\titleclass{\chapter}{straight}
\titleformat{\chapter}[hang]{\normalsize\bfseries}{%
  \thechapter.}{0.5em}{}
\titlespacing*{\chapter}{0pt}{1.5ex plus 0.5ex minus 0.2ex}{0.8ex plus 0.2ex}

\titleformat{\section}[hang]{\normalsize\bfseries}{%
  \thesection}{0.5em}{}
\titlespacing*{\section}{0pt}{1.0ex plus 0.3ex minus 0.2ex}{0.5ex plus 0.1ex}

\titleformat{\subsection}[hang]{\normalsize\bfseries}{%
  \thesubsection}{0.5em}{}
\titlespacing*{\subsection}{0pt}{0.8ex plus 0.3ex minus 0.1ex}{0.4ex plus 0.1ex}

% ── Theorem environments ───────────────────────────────────────────────────────
\newtheorem{theorem}{Theorem}[section]
\newtheorem{lemma}[theorem]{Lemma}
\newtheorem{proposition}[theorem]{Proposition}
\newtheorem{corollary}[theorem]{Corollary}
\newtheorem{definition}[theorem]{Definition}
\newtheorem{example}[theorem]{Example}
\newtheorem{remark}[theorem]{Remark}
\newtheorem{conjecture}[theorem]{Conjecture}
\newtheorem{observation}[theorem]{Observation}

% ── LeanBlueprint macros ──────────────────────────────────────────────────────
% These macros annotate theorem-like environments with Lean 4 declaration
% linkages.  They are consumed by the proof-objects extraction pipeline
% and (optionally) by leanblueprint/plasTeX for dependency graph generation.
%
% Usage:
%   \begin{theorem}\label{thm:my-result}
%     \lean{QOU.Chapter.my_result}      % Lean declaration name
%     \leanok                            % proof is sorry-free
%     \uses{def:prerequisite}            % dependency labels
%     Statement here.
%   \end{theorem}
% Base URL for Lean documentation (doc-gen4 pages on GitHub Pages)
\newcommand{\leandocsbase}{https://litlfred.github.io/qou/lean/}
% GitHub source URL for Lean files
\newcommand{\leansrcbase}{https://github.com/litlfred/qou/blob/main/lean/}
%
% Internal flag: tracks whether \lean was called in the current environment.
\newif\ifleantagged
%
% \lean{Decl.Name} — renders a clickable ∀ icon linking to the Lean declaration.
% Colour depends on formalization status (set by \leanok / default).
% Default: goldenrod (stated); override to forest green with \leanok.
\newcommand{\lean}[1]{%
  \global\leantaggedtrue%
  \texorpdfstring{%
    \href{\leandocsbase#1}{%
      \textsuperscript{\textcolor{BurntOrange}{$\forall$}}%
    }%
  }{(Lean: #1)}%
}
%
% \leanok — marks environment as fully formalized (sorry-free).
% Renders a green checkmark.
\newcommand{\leanok}{%
  \texorpdfstring{%
    \textsuperscript{\textcolor{ForestGreen}{\checkmark}}%
  }{[proved]}%
}
%
% \notready — marks environment as not ready to formalize.
% Renders a red cross.
\newcommand{\notready}{%
  \texorpdfstring{%
    \textsuperscript{\textcolor{BrickRed}{$\times$}}%
  }{[not ready]}%
}
%
% \mathlibok — marks declarations merged into mathlib.
% Renders a dark-green double checkmark.
\newcommand{\mathlibok}{%
  \texorpdfstring{%
    \textsuperscript{\textcolor{OliveGreen}{\checkmark\kern-0.3em\checkmark}}%
  }{[mathlib]}%
}
\newcommand{\uses}[1]{}            % \uses{label1, label2, ...}
\newcommand{\proves}[1]{}          % \proves{thm:label} (in proof environments)
%
% ── Proof-status legend (printed once in the document) ───────────────────────
\newcommand{\proofstatuslegend}{%
  \bigskip\noindent\textbf{Proof-status key}\par\smallskip\noindent
  \begin{tabular}{@{}cl@{\qquad}cl@{}}
    \textcolor{ForestGreen}{\checkmark} & Fully proved (sorry-free)
    & \textcolor{BurntOrange}{$\forall$} & Stated in Lean (has \texttt{sorry})\\
    \textcolor{OliveGreen}{\checkmark\kern-0.3em\checkmark} & Merged into Mathlib
    & \textcolor{BrickRed}{$\times$} & Not ready to formalize\\
  \end{tabular}\bigskip
}

% ── Print mode: affiliations ──────────────────────────────────────────────────
% In compact mode (default), affiliations are omitted.
% Set \showaffiliationstrue for formal print view.
% The MCP render tool writes print-mode.tex to set this flag automatically.
\newif\ifshowaffiliations
\IfFileExists{print-mode.tex}{\input{print-mode.tex}}{}
